\section{上下文管理：最容易被忽视的关键}

Claude Code 的上下文窗口最大可达 1M token，但它不是无限的，且上下文质量直接决定输出质量。

\subsection{上下文里都装了什么}

\begin{itemize}[leftmargin=1.5em]
  \item 对话历史
  \item Claude 读取的所有文件内容
  \item 命令执行的输出
  \item \texttt{CLAUDE.md} 文件（每次都加载）
  \item Memory 文件（前 200 行）
  \item 加载的 Skill 和 MCP 工具定义
\end{itemize}

\subsection{五个实用策略}

\begin{importantbox}{上下文管理五板斧}
\begin{enumerate}[leftmargin=1.5em]
  \item \textbf{\texttt{/clear}}——一件事做完就清。不要在同一个对话里又修 bug 又加功能又重构。
  \item \textbf{\texttt{/compact}}——手动压缩上下文。更好的用法是给压缩加焦点：\texttt{/compact 专注于 API 改动和测试结果}。
  \item \textbf{\texttt{/context}}——查看上下文使用情况，发现哪些 MCP server 占了大量空间。
  \item \textbf{用子 Agent 隔离噪声}——跑测试、分析日志等产生大量输出的操作交给子 Agent，主对话保持干净。
  \item \textbf{在 CLAUDE.md 里写压缩保留指令}——告诉 Claude 压缩时必须保留什么关键信息。
\end{enumerate}
\end{importantbox}

\begin{warningbox}{Memory 的适用边界}
Memory 适合存：偏好、项目约定、历史决策。\\
Memory 不适合存：代码细节、Git 历史、临时状态——这些从代码和 Git 里获取更准确。\\
\texttt{CLAUDE.md} 是团队共享的规范；Memory 是个人的记忆。
\end{warningbox}

\subsection{本章小结}

上下文是稀缺资源。\texttt{/clear} 保持对话聚焦，\texttt{/compact} 延长对话寿命，子 Agent 隔离噪声输出。

%% ====================================================================
